For a database $\mathcal X\subset\mathbb R^d$ and query $q$, the target is
the top-$k$ set under raw squared L2. The implemented system routes through
IVF, ranks the visited vectors by a primary reconstructed score
$\primarydist$, retains up to $c=\max\{k,\lceil\alpha k\rceil\}$ candidates,
and selects results using a residual-refined score $\refineddist$.
Both scores differ from raw-vector distance. The refined value additionally
contains a query-only constant, as defined in \cref{sec:2-3}.
Exactness of a selection procedure therefore concerns its computed score
and tie contract; final ANN quality remains an empirical question against GT.

The representation follows JHQ: a square orthogonal rotation, Cartesian
primary codes and residuals relative to primary reconstruction. This thesis
studies GPU training/encoding, physical placement, query execution and one
workload-level refinement multiplier. The multiplier is selected from
unlabelled query outputs against a larger-budget reference, then is intended
to be reused under a stable workload. It is not a per-query learned controller
or an absolute-recall certificate.

The main cross-system experiments use one RTX 5090, squared-L2 top-10 and
batches of 1,000 queries on four datasets. Six datasets support component
and budget studies, with separate large-index measurements. T5 adds
cross-system batches 8/32 alongside internal scheduling ablations
(\cref{sec:7-smallbatch}). The implemented compressed index
resides on one GPU; CPU model preparation and host-query transfers remain
part of its actual architecture. Multi-GPU distribution, online updates,
host-backed residual paging and billion-scale deployment are outside the
tested system.
